%% Annotated bibliography / literature companion for "Babel's Machines"
%% Each entry: bib key — full citation — verifiable source link — what it does — how Babel's Machines uses it
%% All entries verified against arXiv / ACL Anthology / journal websites (audit: docs/reviews/citation_audit_round1.md)
%% Round 1 + Round 2 reviewer additions also web-verified at insertion time.

\documentclass[11pt,a4paper]{article}
\usepackage[T5,T1]{fontenc}
\usepackage[utf8]{inputenc}
\usepackage[margin=2.5cm]{geometry}
\usepackage[colorlinks=true,linkcolor=blue!60!black,urlcolor=blue!60!black,citecolor=blue!60!black]{hyperref}
\usepackage{xcolor}
\usepackage{enumitem}
\usepackage{titlesec}
\usepackage{parskip}

\definecolor{rsnavy}{HTML}{1B2A4A}

\titleformat{\section}{\Large\bfseries\color{rsnavy}}{\thesection}{0.6em}{}
\titleformat{\subsection}{\normalsize\bfseries}{}{0em}{}

\setlist[itemize]{leftmargin=*,itemsep=0.2em,topsep=0.3em}

% Compact entry macro
% \entry{bibkey}{Citation}{verification URL}{What it does}{How we use it}
\newcommand{\entry}[5]{%
  \par\smallskip
  \noindent\textbf{\texttt{#1}} --- #2\\
  {\footnotesize \textcolor{gray}{Verify:} \url{#3}}\\
  \textit{What it is.} #4\\
  \textit{How Babel's Machines uses it.} #5
  \par\medskip
}

\title{\textbf{Annotated Literature Companion}\\\large for ``Babel's Machines: Linguistic Relativity, Cooperation, and the Emerging Behaviour of Multilingual Artificial Agents''}
\author{Compiled: 2026-04-27}
\date{}

\begin{document}
\maketitle

\section*{Foreword}

This companion document explains every paper cited in \emph{Babel's Machines}. For each entry we provide: the bib key, full citation, a \textbf{verifiable source URL}, a one-paragraph description of what the paper does, and one paragraph on how the present perspective uses, extends, contradicts, or builds on it.

\textbf{Verification.} All cited entries were independently verified against arXiv, ACL Anthology, journal websites (Nature, Science, PNAS, Royal Society, Cambridge Core, IOS Press, MIT Press), and OpenReview. The full audit trail is in \texttt{docs/reviews/citation\_audit\_round1.md}. Reviewer Round 1 and Round 2 additions were web-verified at insertion time. Three categories were distinguished:

\begin{itemize}
  \item \textbf{Verified clean (60+ entries):} title, authors, year, venue, and arXiv ID match authoritative sources.
  \item \textbf{Fixed during audit (12 entries):} discrepancies (wrong venue, missing authors, wrong title, wrong year) were corrected. The most consequential were \texttt{han2022pledges} (venue was incorrectly listed as \emph{Royal Society Open Science}; corrected to ALIFE 2022) and \texttt{turner2023actadd} (wrong title and missing author).
  \item \textbf{Cannot fully verify (1 entry):} \texttt{repello2024multilingualguardrails} is a blog post; the title and content match \url{https://repello.ai/blog/}, but date metadata could not be independently confirmed.
\end{itemize}

Entries are grouped thematically, not alphabetically, so the reader can see the structure of the literature the paper draws on.

\tableofcontents

% =====================================================
\section{Foundations: classical linguistic relativity and the cognitive baseline}

\entry{boroditsky2001does}
{Boroditsky, L.\ (2001). Does language shape thought? Mandarin and English speakers' conceptions of time. \emph{Cognitive Psychology} 43(1):1--22.}
{https://doi.org/10.1006/cogp.2001.0748}
{Classical experimental study showing that Mandarin and English speakers conceptualise time on different spatial axes (vertical vs.\ horizontal), depending on the metaphors built into their respective languages.}
{We cite this as the modern, measured form of the linguistic-relativity hypothesis. The Algorithmic Sapir--Whorf Switch is named after this tradition; we explicitly distance ourselves from the strong determinist Whorfian reading and adopt Boroditsky's ``cognitive affordances'' framing.}

\entry{boroditsky2011languageshapes}
{Boroditsky, L.\ (2011). How language shapes thought. \emph{Scientific American} 304(2):62--65.}
{https://doi.org/10.1038/scientificamerican0211-62}
{Popular-science synthesis of two decades of evidence for measured linguistic relativity, covering colour categorisation, spatial orientation in Kuuk Thaayorre, and grammatical gender effects.}
{Used as the second pillar of our linguistic-relativity background. Together with Boroditsky 2001, it grounds our claim that ``language as cognitive affordance'' is a credible framework to translate to machine populations.}

% =====================================================
\section{The machine-behaviour programme}

\entry{rahwan2019machinebehaviour}
{Rahwan, I., Cebrian, M., Obradovich, N., et al.\ (2019). Machine behaviour. \emph{Nature} 568(7753):477--486.}
{https://doi.org/10.1038/s41586-019-1138-y}
{Manifesto-style \emph{Nature} paper proposing that AI systems be studied as a class of behaving agents using tools from biology, economics, and social science --- not solely the engineering discipline that built them.}
{This is the anchor of our entire argument. The opening of the paper takes Rahwan's programme as the point of departure; we frame the perspective as a contribution within it. Iyad Rahwan is also a guest editor of the theme issue.}

\entry{han2022emergent}
{Han, T.A.\ (2022). Emergent behaviours in multi-agent systems with evolutionary game theory. \emph{AI Communications} 35(4):327--337.}
{https://doi.org/10.3233/AIC-220104}
{Argues that evolutionary game theory provides the natural lens for studying emergent behaviour in multi-agent systems, including those whose strategies arise from learning rather than mutation.}
{We cite this as the methodological bridge between machine populations and evolutionary game theory. Han is a guest editor of the theme issue; the citation is intentional and substantive.}

\entry{han2022pledges}
{Han, T.A., Santos, F.C., Pereira, L.M., Lenaerts, T.\ (2022). Voluntary safety pledges overcome over-regulation dilemma in AI development. \emph{Proc.\ ALIFE 2022}.}
{https://doi.org/10.1162/isal_a_00484}
{Models AI-development regulation as an evolutionary game; shows voluntary pledges with peer or institutional sanctions can resolve the over-regulation dilemma.}
{Cited in our governance discussion as the formal precedent for treating multilingual-AI deployment regulation as an evolutionary-game-theoretic problem. \textbf{Note:} the original bib entry incorrectly listed \emph{Royal Society Open Science} as the venue; corrected to ALIFE 2022 during the audit.}

% =====================================================
\section{Cooperation: the evolutionary game-theory canon}

\entry{axelrod1984evolution}
{Axelrod, R.\ (1984). \emph{The Evolution of Cooperation}. Basic Books, New York.}
{https://www.basicbooks.com/titles/robert-axelrod/the-evolution-of-cooperation/9780465005642/}
{The foundational book demonstrating, via computer tournaments, that cooperative strategies (Tit-for-Tat) can emerge and stabilise among self-interested agents through repeated interaction.}
{We use Axelrod as the negative anchor for the ``culturally inherited cooperation'' argument: the Axelrodian view says cooperation is achieved through repeated interaction and reputation; LLMs cooperate \emph{before} any of those scaffolds operate.}

\entry{nowak2006five}
{Nowak, M.A.\ (2006). Five rules for the evolution of cooperation. \emph{Science} 314(5805):1560--1563.}
{https://doi.org/10.1126/science.1133755}
{Synthesises the five mechanisms by which cooperation can evolve: kin selection, direct reciprocity, indirect reciprocity, network reciprocity, group selection.}
{Cited as the canonical statement of the mechanisms that LLM cooperation does \emph{not} require, supporting our claim that LLM prosocial defaults are a different kind of phenomenon entirely --- cultural inheritance, not evolved reciprocity.}

\entry{nowak1998evolution}
{Nowak, M.A., Sigmund, K.\ (1998). Evolution of indirect reciprocity by image scoring. \emph{Nature} 393(6685):573--577.}
{https://doi.org/10.1038/31225}
{Establishes that reputation-based ``image scoring'' enables cooperation to evolve when direct reciprocity fails.}
{One of the canonical scaffolds whose absence is consequential for the LLM case: LLMs cooperate without reputation, image scoring, or repeated interaction. The contrast sharpens the cultural-inheritance reading.}

\entry{panchanathan2004indirect}
{Panchanathan, K., Boyd, R.\ (2004). Indirect reciprocity can stabilize cooperation without the second-order free rider problem. \emph{Nature} 432(7016):499--502.}
{https://doi.org/10.1038/nature02978}
{Shows that indirect reciprocity, properly modelled, escapes the second-order free-rider problem and stabilises cooperation in larger populations.}
{Cited in the evolutionary section as another scaffold whose absence in LLM populations the cultural-inheritance reading must explain.}

\entry{perc2017statistical}
{Perc, M., Jordan, J.J., Rand, D.G., Wang, Z., Boccaletti, S., Szolnoki, A.\ (2017). Statistical physics of human cooperation. \emph{Physics Reports} 687:1--51.}
{https://doi.org/10.1016/j.physrep.2017.05.004}
{Comprehensive review of the statistical-physics approach to cooperation: spatial reciprocity, network reciprocity, phase transitions, and the dynamics of cooperative regimes.}
{Cited as the second pillar of the theoretical machinery we propose to apply to LLM populations. Mat\v{j}az Perc is a guest editor of the theme issue; the link to his programme is intentional.}

\entry{neumann1944theory}
{von Neumann, J., Morgenstern, O.\ (1944). \emph{Theory of Games and Economic Behavior}. Princeton University Press.}
{https://press.princeton.edu/books/paperback/9780691130613/theory-of-games-and-economic-behavior}
{The book that founded modern game theory as a discipline.}
{Foundational acknowledgement only; not centrally argued from.}

\entry{nash1950equilibrium}
{Nash, J.F.\ (1950). Equilibrium points in $n$-person games. \emph{Proc.\ National Academy of Sciences} 36(1):48--49.}
{https://doi.org/10.1073/pnas.36.1.48}
{The two-page note that introduced the Nash equilibrium concept.}
{Foundational acknowledgement only.}

% =====================================================
\section{Cultural evolution and the WEIRD problem}

\entry{boyd1985culture}
{Boyd, R., Richerson, P.J.\ (1985). \emph{Culture and the Evolutionary Process}. University of Chicago Press.}
{https://press.uchicago.edu/ucp/books/book/chicago/C/bo3621168.html}
{Foundational work formalising cultural evolution as a Darwinian process operating on transmitted information, with mechanisms (oblique/horizontal/vertical inheritance, conformist transmission, etc.) analogous to genetic evolution.}
{Central to our reframe. The claim that pretraining + RLHF act as cultural inheritance is explicitly framed in Boyd--Richerson terms; LLMs are cast as a new substrate for their framework.}

\entry{richerson2005notbygenes}
{Richerson, P.J., Boyd, R.\ (2005). \emph{Not by Genes Alone: How Culture Transformed Human Evolution}. University of Chicago Press.}
{https://press.uchicago.edu/ucp/books/book/chicago/N/bo3636224.html}
{Companion volume to Boyd \& Richerson 1985, accessible to a wider audience, summarising the case for cumulative cultural evolution.}
{Cited together with Boyd 1985 to anchor the cultural-inheritance reading.}

\entry{henrich2016secret}
{Henrich, J.\ (2015). \emph{The Secret of Our Success: How Culture Is Driving Human Evolution\dots}. Princeton University Press.}
{https://press.princeton.edu/books/hardcover/9780691166858/the-secret-of-our-success}
{Argues that humans are an ultra-cooperative species because of cumulative culture, not genes; cooperation is built up by transmission, not selected for biologically.}
{Used as the third anchor of the cultural-inheritance argument. Henrich's transmission framing is what we apply to machine-text-trained agents.}

\entry{henrich2010weirdest}
{Henrich, J., Heine, S.J., Norenzayan, A.\ (2010). The weirdest people in the world? \emph{Behavioral and Brain Sciences} 33(2--3):61--83.}
{https://doi.org/10.1017/S0140525X0999152X}
{Shows that Western, Educated, Industrialised, Rich, Democratic (WEIRD) populations are statistical outliers on many traits the cooperation literature treats as universal.}
{This is the linchpin of the WEIRD-corollary argument. If LLM cooperation is culturally inherited from WEIRD-skewed corpora, LLMs encode WEIRD norms while presenting as universal --- a normative externality of multilingual deployment.}

\entry{muthukrishna2020beyondweird}
{Muthukrishna, M., Bell, A.V., Henrich, J., Curtin, C.M., Gedranovich, A., McInerney, J., Thue, B.\ (2020). Beyond WEIRD Psychology. \emph{Psychological Science} 31(6):678--701.}
{https://doi.org/10.1177/0956797620916782}
{Operationalises ``WEIRD'' as a measurable cultural distance and maps where in the world WEIRDness deviates from local norms.}
{Cited together with Henrich 2010 to make the WEIRD argument quantitative rather than rhetorical. Provides the methodology to test whether LLM defaults track WEIRD distance.}

\entry{tomasello2009whywecooperate}
{Tomasello, M.\ (2009). \emph{Why We Cooperate}. MIT Press.}
{https://mitpress.mit.edu/9780262013598/why-we-cooperate/}
{Developmental and comparative argument that human cooperation rests on shared intentionality and joint attention --- cognitive primitives largely absent in other apes.}
{Cited as part of the broader cooperation-and-cognition canon that the perspective draws on; supports the claim that ToM is a cognitive prerequisite for cooperation.}

% =====================================================
\section{Theory of mind: foundations and modern critique}

\entry{premack1978tom}
{Premack, D., Woodruff, G.\ (1978). Does the chimpanzee have a theory of mind? \emph{Behavioral and Brain Sciences} 1(4):515--526.}
{https://doi.org/10.1017/S0140525X00076512}
{The 1978 paper that named ``theory of mind'' as a research target in comparative cognition.}
{Cited as the etymological anchor when we introduce ToM in Section 5.}

\entry{deweerd2013bayesiantom}
{de Weerd, H., Verbrugge, R., Verheij, B.\ (2013). How much does it help to know what she knows you know? \emph{Artificial Intelligence} 199--200:67--92.}
{https://doi.org/10.1016/j.artint.2013.05.004}
{Agent-based simulation establishing that recursive ToM (k-level reasoning) provides diminishing returns past first or second order, with first-order being the strongest predictor of cooperative success.}
{Cited as the formal ground for the LLM-Hanabi finding that \emph{first-order} ToM is what predicts cooperation in LLM agents --- consistent with the human/AI agent literature.}

\entry{byrne1988machiavellian}
{Byrne, R.W., Whiten, A.\ (eds.) (1988). \emph{Machiavellian Intelligence}. Oxford University Press.}
{https://global.oup.com/academic/product/machiavellian-intelligence-9780198521754}
{The volume that established the ``Machiavellian intelligence'' hypothesis: that primate cognitive evolution was driven by the demands of social manipulation, not domain-general problem-solving.}
{Cited in the SPIN-Bench discussion: the planning-vs-social-intelligence dissociation in LLMs mirrors the evolutionary distinction between domain-general planning and Machiavellian social cognition.}

\entry{ullman2023tomfail}
{Ullman, T.\ (2023). Large Language Models Fail on Trivial Alterations to Theory-of-Mind Tasks. arXiv:2302.08399.}
{https://arxiv.org/abs/2302.08399}
{Demonstrates that small variations to standard false-belief items collapse apparent LLM ToM performance, suggesting reported ToM successes may be surface pattern-matching rather than genuine mental-state representation.}
{Round 2 reviewer addition. Cited explicitly in the ToM construct-validity caveat: even if absolute LLM ToM scores are contestable, our claim about \emph{cross-lingual differences} on the same items remains well-defined.}

\entry{sap2022neuraltom}
{Sap, M., LeBras, R., Fried, D., Choi, Y.\ (2022). Neural Theory-of-Mind? On the Limits of Social Intelligence in Large LMs. \emph{EMNLP 2022}. arXiv:2210.13312.}
{https://arxiv.org/abs/2210.13312}
{Reports that GPT-3, ChatGPT, and GPT-4 reach only $\sim 60\%$ on SocialIQa and ToMi benchmarks, well below human; argues scaling alone is insufficient for ToM.}
{Round 2 reviewer addition. Pairs with Ullman 2023 in the construct-validity caveat. The combined citation acknowledges the controversy without conceding the cross-lingual asymmetry argument.}

\entry{kovacs2009bilingual}
{Kov\'acs, \'A.M.\ (2009). Early bilingualism enhances mechanisms of false-belief reasoning. \emph{Developmental Science} 12(1):48--54.}
{https://doi.org/10.1111/j.1467-7687.2008.00742.x}
{Shows that bilingual children pass false-belief tasks earlier than monolingual peers, linking the executive demands of code-switching to ToM.}
{This is the human counter-prediction that frames the LLM ToM result as surprising: in humans, multilingualism \emph{strengthens} ToM; in LLMs, it weakens it. The contrast is rhetorical and substantive.}

\entry{goetz2003bilingual}
{Goetz, P.J.\ (2003). The effects of bilingualism on theory of mind development. \emph{Bilingualism: Language and Cognition} 6(1):1--15.}
{https://doi.org/10.1017/S1366728903001007}
{Independent confirmation of the bilingual ToM advantage with a different developmental population.}
{Paired with Kov\'acs 2009 to make the bilingual-ToM finding robust rather than a one-study artefact.}

% =====================================================
\section{Foundational LLM architecture and training}

\entry{vaswani2017transformer}
{Vaswani, A., Shazeer, N., Parmar, N., et al.\ (2017). Attention is all you need. \emph{NeurIPS 2017}.}
{https://arxiv.org/abs/1706.03762}
{The original Transformer architecture paper; introduced self-attention as the dominant primitive for sequence modelling.}
{Cited as the architectural substrate; everything we discuss assumes Transformer-based LLMs.}

\entry{brown2020gpt3}
{Brown, T., Mann, B., Ryder, N., et al.\ (2020). Language models are few-shot learners. \emph{NeurIPS 2020}.}
{https://arxiv.org/abs/2005.14165}
{The GPT-3 paper; showed few-shot in-context learning emerges at scale and that English-dominant pretraining produces an asymmetric multilingual model.}
{Cited when we note that pretraining corpora are largely English-dominant; this is the statistical asymmetry that the Semantic Hub Hypothesis later explains.}

\entry{ouyang2022training}
{Ouyang, L., Wu, J., Jiang, X., et al.\ (2022). Training language models to follow instructions with human feedback. \emph{NeurIPS 2022}.}
{https://arxiv.org/abs/2203.02155}
{The InstructGPT paper; established RLHF as a standard alignment recipe.}
{Cited where we discuss RLHF as the second leg of cultural inheritance and as a documented source of linguistic skew (RLHF annotators are disproportionately English-speaking).}

\entry{bai2022constitutional}
{Bai, Y., Kadavath, S., Kundu, S., et al.\ (2022). Constitutional AI: Harmlessness from AI feedback. arXiv:2212.08073.}
{https://arxiv.org/abs/2212.08073}
{Anthropic's RLAIF approach: replaces human preference labels with AI-generated critiques against a written ``constitution''.}
{Cited alongside Ouyang to broaden the RLHF discussion to include AI-feedback variants; the linguistic-skew problem persists in RLAIF if the critic model is itself English-skewed.}

% =====================================================
\section{Multilingual representation: classics}

\entry{conneau2020xlmr}
{Conneau, A., Khandelwal, K., Goyal, N., et al.\ (2020). Unsupervised cross-lingual representation learning at scale (XLM-R). \emph{ACL 2020}.}
{https://arxiv.org/abs/1911.02116}
{Trained a multilingual masked language model on 100 languages and 2.5TB of CommonCrawl; established that scale + balanced sampling can rival monolingual models on individual languages.}
{Round 1 reviewer addition. Cited in the architectural-alternatives subsection (\S 4.1): scale alleviates but does not eliminate cross-lingual subspace interference.}

\entry{xue2021mt5}
{Xue, L., Constant, N., Roberts, A., et al.\ (2021). mT5: A massively multilingual pre-trained text-to-text transformer. \emph{NAACL 2021}.}
{https://arxiv.org/abs/2010.11934}
{Multilingual T5 across 101 languages; documented ``accidental translation'' (the model producing output in unintended languages) as a failure mode.}
{Round 1 reviewer addition. Cited together with XLM-R as evidence that subspace interference persists even in well-engineered multilingual models. Also cited as an example of a calibration-repair lever (language adapters / instruction-tuning variations).}

\entry{pires2019multilingualbert}
{Pires, T., Schlinger, E., Garrette, D.\ (2019). How multilingual is multilingual BERT? \emph{ACL 2019}.}
{https://arxiv.org/abs/1906.01502}
{Showed that multilingual BERT supports surprising zero-shot cross-lingual transfer but exhibits ``systematic deficiencies'' on certain language pairs.}
{Round 1 reviewer addition. The ``systematic deficiencies'' finding is the empirical base for the language-specific microcircuits alternative to the Semantic Hub.}

% =====================================================
\section{Mechanistic interpretability: hub, barrier, and alternatives}

\entry{wu2024semantichub}
{Wu, Z., Yu, X.V., Yogatama, D., Lu, J., Kim, Y.\ (2024/2025). The Semantic Hub Hypothesis. \emph{ICLR 2025}.}
{https://arxiv.org/abs/2411.04986}
{Provides causal evidence (via activation steering) that multilingual LLMs project heterogeneous linguistic inputs into a centralised, English-dominant intermediate hub.}
{This is one of the two architectural pillars of Section 4. We cite it for both the representational claim and the activation-steering methodology that we propose to extend with a behavioural readout in P2.}

\entry{wendler2024llamasenglish}
{Wendler, C., Veselovsky, V., Monea, G., West, R.\ (2024). Do Llamas Work in English? On the Latent Language of Multilingual Transformers. \emph{ACL 2024}.}
{https://aclanthology.org/2024.acl-long.820/}
{Logit-lens evidence that Llama-family decoders pass through an intermediate state that is statistically closer to English than to the input language.}
{Cited as the ``latent-language priors'' alternative to the Semantic Hub --- a stronger claim that the hub is specifically English-anchored, not just shared. The natural target for the depth-resolved probes proposed in P2.}

\entry{bafna2025translationbarrier}
{Bafna, N., Li, T., Murray, K., Mortensen, D.R., Yarowsky, D., Sirin, H., Khashabi, D.\ (2025). The Translation Barrier Hypothesis. \emph{IJCNLP-AACL 2025}.}
{https://arxiv.org/abs/2506.22724}
{Shows that many multilingual generation failures are not upstream-reasoning failures but late-layer realisation failures: the model reaches the right conclusion internally and then fails projecting into the target language.}
{The second architectural pillar. Combined with the Semantic Hub, gives the hub-vs-realisation diagnostic categories that P2 operationalises.}

\entry{turner2023actadd}
{Turner, A.M., Thiergart, L., Leech, G., Udell, D., Vazquez, J.J., Mini, U., MacDiarmid, M.\ (2023). Steering Language Models With Activation Engineering. arXiv:2308.10248.}
{https://arxiv.org/abs/2308.10248}
{Introduced Activation Addition (ActAdd) as a lightweight, training-free way to steer LLM behaviour by adding activation differences at inference time.}
{Methodological precedent for the steering-vector experiment in P2. \textbf{Note:} the original bib entry had the wrong title (``Activation Addition: Steering Language Models Without Optimization'') and a missing author (Vazquez); both fixed during the audit.}

\entry{kumar2025indictunedlens}
{Panchal, M., Varshney, D., Mamta, Ekbal, A.\ (2025/2026). Indic-TunedLens: Interpreting Multilingual Models in Indian Languages. \emph{EACL 2026 VarDial Workshop}.}
{https://arxiv.org/abs/2602.15038}
{Adapts the TunedLens technique to Indic languages, providing per-language layer-resolved interpretability tooling for low-resource scripts.}
{Cited as the methodology for distinguishing hub failures from realisation failures in P2; without per-language probes, the hub/barrier diagnostic cannot be operationalised in low-resource languages.}

% =====================================================
\section{Interpretability: causal-evidence precedents (Round 2 additions)}

\entry{chen2023spatialcausal}
{Chen, Y., Gan, Y., Li, S., Yao, L., Zhao, X.\ (2023). More than Correlation: Do Large Language Models Learn Causal Representations of Space? arXiv:2312.16257.}
{https://arxiv.org/abs/2312.16257}
{Causal-intervention experiments showing that spatial representations in DeBERTa and GPT-Neo \emph{causally} affect next-token prediction, not merely correlate with it.}
{Round 2 addition. Cited in \S 4.3 (Closing the Causal Loop) as the methodological precedent for the steering-vector $\to$ behaviour experiment in P2.}

\entry{zhang2026worldmodels}
{Zhang, X.\ (2026). What Do World Models Learn in RL? Probing Latent Representations in Learned Environment Simulators. arXiv:2603.21546.}
{https://arxiv.org/abs/2603.21546}
{Demonstrates that latent state variables in IRIS and DIAMOND world models are linearly decodable \emph{and} causally engaged via targeted intervention.}
{Round 2 addition. Used as a second causal-intervention precedent that justifies the move from representational diagnosis to behavioural readout in P2.}

\entry{yan2025arithmetic}
{Yan, Y.\ (2025). Addition in Four Movements: Mapping Layer-wise Information Trajectories in LLMs. arXiv:2506.07824.}
{https://arxiv.org/abs/2506.07824}
{Identifies a coherent four-stage layer-by-layer trajectory in Llama-3 doing multi-digit addition: formula representation $\to$ numerical features $\to$ digit-by-digit results $\to$ final token.}
{Round 2 addition. Cited as evidence that depth-resolved probing can recover meaningful processing stages, justifying the L/3 and 2L/3 probe choices in P2.}

\entry{lee2024causalpaths}
{Lee, I., Lum, J., Liu, Z., Yogatama, D.\ (2024). Causal Interventions on Causal Paths: Mapping GPT-2's Reasoning From Syntax to Semantics. arXiv:2410.21353.}
{https://arxiv.org/abs/2410.21353}
{Causal intervention study showing GPT-2 processes causal-language sentences via early-layer syntactic detection followed by late-layer semantic analysis.}
{Round 2 addition. Cited as direct precedent for the early-layers-vs-late-layers dissociation that P2's hub-vs-realisation split assumes.}

\entry{zhenyu2025logitlens}
{Hong, Z.\ (2025). LogitLens4LLMs: Extending Logit Lens Analysis to Modern Large Language Models. arXiv:2503.11667.}
{https://arxiv.org/abs/2503.11667}
{Open-source toolkit extending the Logit Lens technique to Llama-3.1, Qwen-2.5, and other modern architectures, with memory-efficient hooks for layer-wise analysis up to 70B parameters.}
{Round 2 addition. Cited as the practical tooling that makes P2 immediately runnable on frontier models.}

% =====================================================
\section{Calibration across languages}

\entry{ahuja2022calibration}
{Ahuja, K., Sitaram, S., Dandapat, S., Choudhury, M.\ (2022). On the Calibration of Massively Multilingual Language Models. \emph{EMNLP 2022}.}
{https://aclanthology.org/2022.emnlp-main.290/}
{Documents the cross-lingual calibration gap on the MEGA benchmark across XGLM, BLOOM, and PaLM with 15-bin equal-width binning; reports ECE rising from $\sim 4.6\%$ in English to $\sim 23.1\%$ in non-English.}
{The empirical source for the ``$\sim 5\times$ ECE gap'' claim that anchors the calibration-crisis discussion in Section 4.}

\entry{zhou2025beyond}
{Zhou, E., Zhang, C., Hu, T., Li, C., Collier, N., Vuli\'c, I., Korhonen, A.\ (2025). Beyond the Final Layer: Intermediate Representations for Better Multilingual Calibration. arXiv:2510.03136.}
{https://arxiv.org/abs/2510.03136}
{Reproduces the cross-lingual ECE gap on Llama-3 and Qwen-2 and shows that probing a late-intermediate layer (rather than the final layer) recovers part of the gap.}
{Cited together with Ahuja 2022 to show the ECE asymmetry replicates across model generations and to motivate intermediate-layer probing in P2.}

\entry{fu2025multilingualjudge}
{Fu, X., Liu, W.\ (2025). How Reliable is Multilingual LLM-as-a-Judge? \emph{Findings of EMNLP 2025}.}
{https://aclanthology.org/2025.findings-emnlp.587/}
{Empirically establishes that single-judge LLM-as-Judge reliability collapses to Fleiss' $\kappa \approx 0.3$ across languages.}
{Direct empirical anchor for P4 (Evaluation reliability ceiling). The whole prediction is structured around the question of whether ensemble protocols can lift this ceiling above $0.5$.}

% =====================================================
\section{Multilingual safety, jailbreaks, and adversarial probes}

\entry{yong2024lowresource}
{Yong, Z.-X., Menghini, C., Bach, S.H.\ (2023). Low-Resource Languages Jailbreak GPT-4. \emph{NeurIPS 2023 SoLaR Workshop} (Best Paper).}
{https://arxiv.org/abs/2310.02446}
{Demonstrates that translating prompts into low-resource languages (Zulu, Scots Gaelic, Hmong, Guarani) bypasses GPT-4's safety training.}
{Cited as evidence for the ``RLHF safety coverage is skewed'' confound in Section 3.3, and for the broader claim that linguistic asymmetries open attack surfaces.}

\entry{yong2025multilingualsafety}
{Yong, Z.-X., Ermis, B., Fadaee, M., Bach, S.H., Kreutzer, J.\ (2025). The State of Multilingual LLM Safety Research: From Measuring the Language Gap to Mitigating It. \emph{EMNLP 2025}.}
{https://aclanthology.org/2025.emnlp-main.800/}
{Survey of the multilingual safety literature; documents the gap between English and non-English safety coverage and proposes mitigation directions.}
{Cited as the modern survey anchor for the ``RLHF coverage is skewed'' claim and as one of the safety benchmarks to be paired with interactive multi-agent ones in our governance recommendations.}

\entry{yoo2025csrt}
{Yoo, H., Yang, Y., Lee, H.\ (2025). Code-Switching Red-Teaming. \emph{ACL 2025}.}
{https://aclanthology.org/2025.acl-long.657/}
{Shows that intra-utterance language mixing achieves a $46.7\%$ increase in attack success against frontier LLM safety systems.}
{The empirical source for the code-switching attack-surface claim. Cited in Section 7 for societal consequences and again in the governance levers (Code-Switching probes as default audit).}

\entry{he2025agentmiddle}
{He, P., Lin, Y., Dong, S., Xu, H., Xing, Y., Liu, H.\ (2025). Red-Teaming LLM Multi-Agent Systems via Communication Attacks. \emph{Findings of ACL 2025}.}
{https://aclanthology.org/2025.findings-acl.349/}
{Introduces Agent-in-the-Middle (AiTM) attacks that intercept and manipulate inter-agent messages in multilingual multi-agent systems.}
{Cited together with CSRT as the multi-agent extension of the attack-surface argument; both become default audit probes in our governance recommendations.}

\entry{m-alert2025}
{Friedrich, F., Tedeschi, S., Schramowski, P., Brack, M., Navigli, R., Nguyen, H., Li, B., Kersting, K.\ (2025). LLMs Lost in Translation: M-ALERT uncovers Cross-Linguistic Safety Inconsistencies. \emph{ICLR 2025}.}
{https://arxiv.org/abs/2412.15035}
{Multilingual safety benchmark documenting cross-linguistic inconsistency in alignment behaviour.}
{Cited as one of the safety benchmarks to be paired with interactive multi-agent benchmarks in our diversification recommendation. \textbf{Note:} the original bib entry had wrong title and venue (``OpenReview''); both corrected during the audit.}

\entry{wang2024evaluating}
{Huang, K., Mo, F., Zhang, X., Li, H., Li, Y., Zhang, Y., Yi, W., Mao, Y., Liu, J., Xu, Y., Xu, J., Nie, J.-Y., Liu, Y.\ (2024). A Survey on Large Language Models with Multilingualism: Recent Advances and New Frontiers. arXiv:2405.10936.}
{https://arxiv.org/abs/2405.10936}
{Comprehensive survey of multilingual LLM research, covering pretraining, alignment, evaluation, and deployment gaps.}
{Cited as broader survey context in Section 7. \textbf{Note:} earlier bib entry was missing author Xinyu Zhang; corrected during the audit.}

\entry{repello2024multilingualguardrails}
{Repello AI (2024). Introducing New Multilingual AI Safety Guardrails for 100 Languages. Blog post.}
{https://repello.ai/blog/introducing-new-multilingual-ai-safety-guardrails-for-100-languages}
{Industry blog post announcing a multilingual safety guardrail system covering 100 languages.}
{Cited as evidence that the industry is starting to address the multilingual safety gap. \textbf{Note:} this is the one citation the audit could not fully verify (date metadata not confirmable); the user may want to substitute with the academic CREST paper.}

\entry{stanfordhai2024languagegap}
{Stanford HAI Policy Brief (2025). Mind the (Language) Gap: Mapping the Challenges of LLM Development in Low-Resource Language Contexts.}
{https://hai.stanford.edu/policy/mind-the-language-gap}
{Stanford HAI policy brief mapping the structural challenges of LLM development for low-resource languages.}
{Cited in Section 7 as policy-community recognition of the language gap as a structural issue. \textbf{Note:} year corrected from 2024 to 2025 during the audit.}

% =====================================================
\section{Multilingual social-dilemma evidence}

\entry{huynh2026moreatstake}
{Huynh, T.-K., Dao-Sy, D.-M., Cao, T.-B., et al.\ (2026). More at Stake: How Payoff and Language Shape LLM Agent Strategies in Cooperation Dilemmas. arXiv:2601.19082.}
{https://arxiv.org/abs/2601.19082}
{Large-scale multilingual IPD evaluation across four model families and five languages, with strategy-classifier labelling at varying confidence thresholds.}
{The primary empirical source for Table 1. Most of the cross-lingual divergence patterns we discuss in Section 3.1 come from this study (the user's own work).}

\entry{huynh2025fairgame}
{Huynh, T.-K., Dao-Sy, D.-M., Cao, T.-B., et al.\ (2025). Understanding LLM Agent Behaviours via Game Theory. arXiv:2512.07462.}
{https://arxiv.org/abs/2512.07462}
{Companion methodological paper to the IPD study; details strategy-recognition methodology, multi-agent dynamics, and bias analysis.}
{Cited together with the More-at-Stake paper for methodological details, and specifically for the Vietnamese low-stakes ``defection'' = arithmetic-error finding.}

\entry{za2026towards}
{Za, J., Panos, A., Cuhel, J., Albanie, S.\ (2026). Towards Predictive Models of Strategic Behaviour in Large Language Model Agents. \emph{ICLR 2026 Workshop}.}
{https://openreview.net/forum?id=cBzy3OXXm5}
{Builds predictive models of LLM strategic behaviour across IPD variants and reports prosocial bias under multilingual conditions.}
{Used as supporting evidence for the prosocial-default finding and as a multilingual extension of the Fontana ``nicer than human'' result.}

\entry{fontana2024nicer}
{Fontana, N., Pierri, F., Aiello, L.M.\ (2024/2025). Nicer Than Humans: How do Large Language Models Behave in the Prisoner's Dilemma? \emph{ICWSM 2025}.}
{https://arxiv.org/abs/2406.13605}
{Reports that Llama2 and GPT-3.5 are markedly more cooperative than humans in IPD even against random adversaries with substantial defection probabilities.}
{The empirical anchor for the bounded ``nicer than human'' claim in Section 3.2. We rely on this paper for the specific regime (cooperative payoff, disclosed horizon, exploitable opponent) that bounds the prosocial default.}

\entry{wang2026discovering}
{Wang, C., Kasenberg, D., Stachenfeld, K., Castro, P.S.\ (2026). Discovering Differences in Strategic Behavior Between Humans and LLMs. arXiv:2602.10324.}
{https://arxiv.org/abs/2602.10324}
{Shows that in iterated rock-paper-scissors, frontier LLMs play \emph{deeper} strategically than humans --- the inverse of the cooperative-game prosocial bias.}
{Critical for bounding the ``nicer than human'' claim: the prosocial default is genre-specific, not a uniform LLM trait.}

\entry{buscemi2025fairgame}
{Buscemi, A., Proverbio, D., Di Stefano, A., Han, T.A., Castignani, G., Li\`o, P.\ (2025). FAIRGAME: A Framework for AI Agents Bias Recognition using Game Theory. \emph{ECAI 2025} (Outstanding Paper).}
{https://arxiv.org/abs/2504.14325}
{Game-theoretic framework for measuring biases in LLM agents; reports that when game horizons are unknown, cooperation universalises across all languages.}
{Cited as one of the key contradiction findings in Section 3.3: framing (horizon disclosure) can fully dominate language effects.}

\entry{lore2024framing}
{Lor\`e, N., Heydari, B.\ (2024). Strategic Behavior of LLMs and the Role of Game Structure versus Contextual Framing. \emph{Scientific Reports} 14:18490.}
{https://doi.org/10.1038/s41598-024-69032-z}
{Demonstrates that contextual framing of strategic-game prompts can dominate game-theoretic structure in LLM behaviour.}
{Together with FAIRGAME, anchors the ``framing dominates language'' finding in Section 3.3.}

\entry{akata2025playing}
{Akata, E., Schulz, L., et al.\ (2025). Playing Repeated Games with Large Language Models. \emph{Nature Human Behaviour} 9:1380--1390.}
{https://doi.org/10.1038/s41562-025-02172-y}
{High-profile \emph{Nature Human Behaviour} paper systematically evaluating LLM behaviour in repeated games, including coordination and Prisoner's Dilemma.}
{Cited as part of the social-dilemma substrate in Section 2 and as a high-profile baseline for the field.}

% =====================================================
\section{Moderating factors beyond language (Round 1 additions)}

\entry{madmoun2025communication}
{Madmoun, H., Lahlou, S.\ (2025). Communication Enables Cooperation in LLM Agents. arXiv:2510.05748.}
{https://arxiv.org/abs/2510.05748}
{Reports that introducing a one-word communication channel raises Stag-Hunt cooperation from 0\% to 96.7\% across LLM agents, while a defection-biased curriculum induces ``learned pessimism''.}
{Round 1 addition. Cited as the strongest single piece of evidence for the ``cheap talk dominates language'' finding in Section 3.3 and on the contradiction slide.}

\entry{ong2025personalitysteering}
{Ong, K.J.K., Lye, J.J., Nguyen, H.M., Cho, S.H., P\'erez-Campanero Antol\'in, N.\ (2025). Identifying Cooperative Personalities in Multi-agent Contexts through Personality Steering. arXiv:2503.12722.}
{https://arxiv.org/abs/2503.12722}
{Uses representation engineering to steer Big-Five Agreeableness and Conscientiousness in LLMs, modulating IPD cooperation while increasing exploitation susceptibility.}
{Round 1 addition. Cited as direct evidence that internal representations \emph{can} causally steer cooperative play, which strengthens the more specific claim that language-linked representations do so.}

\entry{nishimoto2026commdelay}
{Nishimoto, K., Asatani, K., Sakata, I.\ (2026). Cooperation Breakdown in LLM Agents Under Communication Delays. arXiv:2602.11754.}
{https://arxiv.org/abs/2602.11754}
{Documents a U-shaped relationship between communication delay and cooperation in continuous Prisoner's Dilemma; agents spontaneously begin to exploit slower partners.}
{Round 1 addition. Cited as the infrastructure-modulator example in Section 5.1: behaviour can change without any language change.}

\entry{jiang2025trustcollides}
{Jiang, G., Yang, S., Wang, Y., Hui, P.\ (2025). When Trust Collides: Decoding Human-LLM Cooperation Dynamics. arXiv:2503.07320.}
{https://arxiv.org/abs/2503.07320}
{Shows that human cooperation in a repeated Prisoner's Dilemma against an agent depends on the declared identity of the agent (human / rule-based AI / LLM).}
{Round 1 addition. Cited as the partner-identity-framing modulator in Section 5.1; language is one of several signals that frame agent identity.}

\entry{du2024ctllm}
{Du, X., Yu, Z., Gao, S., et al.\ (2024). Chinese Tiny LLM: Pretraining a Chinese-Centric Large Language Model. arXiv:2404.04167.}
{https://arxiv.org/abs/2404.04167}
{2B-parameter model trained from scratch on 800B Chinese + 300B English + 100B code tokens; documents the training procedure and releases the corpus and benchmarks.}
{Round 1 addition. Cited as one of two language-prioritised models that provide the natural test of whether the Semantic Hub asymmetry is data-driven or architectural.}

\entry{raihan2025tigerllm}
{Raihan, N., Zampieri, M.\ (2025). TigerLLM: A Family of Bangla Large Language Models. arXiv:2503.10995.}
{https://arxiv.org/abs/2503.10995}
{Bangla-centric LLM family that surpasses open-source alternatives and outperforms larger English-centric proprietary models on Bangla benchmarks.}
{Round 1 addition. The Bangla counterpart to CT-LLM in the language-prioritised-models test of the Semantic Hub.}

% =====================================================
\section{Negotiation, persuasion, and auction}

\entry{meta2022cicero}
{Meta FAIR Diplomacy Team (2022). Human-Level Play in the Game of Diplomacy. \emph{Science} 378(6624):1067--1074.}
{https://doi.org/10.1126/science.ade9097}
{The CICERO system: combines language models with strategic planning to play Diplomacy at human level in online tournaments.}
{Cited as the canonical example of LLM agents engaging in long-horizon negotiation. Used in Section 1's catalogue of behaviours and in Section 2's substrate description.}

\entry{deng2024bargaining}
{Deng, Y., Mirrokni, V., et al.\ (2024). LLMs at the Bargaining Table. \emph{ICML 2024 Workshop}.}
{https://openreview.net/}
{Workshop paper analysing LLM behaviour in two-party bargaining settings.}
{Part of the negotiation-and-bargaining family in Section 2.}

\entry{yao2025spinbench}
{Yao, J., Wang, K., Hsieh, R., Zhou, H., Zou, T., Cheng, Z.W., Wang, Z., Viswanath, P.\ (2025). SPIN-Bench: How Well Do LLMs Plan Strategically and Reason Socially? \emph{COLM 2025}.}
{https://arxiv.org/abs/2503.12349}
{Decouples formal classical planning from social intelligence; shows o-series reasoning models lead PDDL planning but degrade chain-of-thought coherence when negotiation phases are introduced.}
{Critical for the planning-vs-ToM dissociation argument in Section 5: cross-lingual transfer is uneven across these faculties and demands separate measurement.}

\entry{biswas2025persuasion}
{Biswas, S., Erlei, A., Gadiraju, U.\ (2025). Mind the Gap! Choice Independence in Using Multilingual LLMs for Persuasive Co-Writing. \emph{CHI 2025}.}
{https://arxiv.org/abs/2502.09532}
{Cross-lingual study of persuasive co-writing for charitable appeals; finds prior LLM-language exposure shifts usage but not aggregate persuasiveness; donor beliefs about AI authorship dominate.}
{The persuasion-null contradiction in Section 3.3. One of the four findings that bounds the linguistic-relativity claim.}

\entry{chen2024aucarena}
{Chen, J., Yuan, S., Ye, R., Majumder, B.P., Richardson, K.\ (2023/2024). AucArena: Evaluating Strategic Planning of LLM Agents in an Auction Arena. arXiv:2310.05746.}
{https://arxiv.org/abs/2310.05746}
{Auction-arena environment for evaluating LLM strategic planning, budget management, and execution.}
{Part of the persuasion-and-auction family in Section 2.}

% =====================================================
\section{Social-deduction games}

\entry{xu2023werewolf}
{Xu, Y., Wang, S., Li, P., Luo, F., Wang, X., Liu, W., Liu, Y.\ (2023). Exploring LLMs for Communication Games: An Empirical Study on Werewolf. arXiv:2309.04658.}
{https://arxiv.org/abs/2309.04658}
{Empirical study of LLM agents playing Werewolf, examining communication, deception, and coalition-formation behaviours.}
{Part of the social-deduction substrate in Section 2. \textbf{Note:} bib entry author list was truncated as ``and others''; expanded to full 7-author list during the audit.}

\entry{wang2023avalon}
{Wang, S., Liu, C., Zheng, Z., Qi, S., Chen, S., Yang, Q., Zhao, A., Wang, C., Song, S., Huang, G.\ (2023). Avalon's Game of Thoughts: Battle Against Deception through Recursive Contemplation. arXiv:2310.01320.}
{https://arxiv.org/abs/2310.01320}
{Avalon-playing LLM agents using recursive contemplation to handle multi-level deception.}
{Part of the social-deduction substrate. \textbf{Note:} bib author list expanded from ``and others'' to full 10-author list during the audit.}

\entry{koto2025spyfall}
{Wibowo, H.A., Elsetohy, A., Cui, Q., Aji, A.F.\ (2026). Multicultural Spyfall: Assessing LLMs through Dynamic Multilingual Social Deduction Game. arXiv:2601.09017.}
{https://arxiv.org/abs/2601.09017}
{Multilingual extension of Spyfall designed to assess cross-cultural social-deduction behaviour in LLMs.}
{Cited as the multilingual social-deduction component of the substrate; relevant precedent for the cross-lingual-divergence argument.}

% =====================================================
\section{Cooperative communication and theory-of-mind games}

\entry{hakimov2025codenames}
{Hakimov, S., Pfennigschmidt, L., Schlangen, D.\ (2025). Ad-hoc Concept Forming in the Game Codenames. \emph{GEM\textsuperscript{2} Workshop @ ACL 2025}.}
{https://aclanthology.org/}
{Uses Codenames as a benchmark to evaluate ad-hoc concept formation and shared cultural reference in LLMs.}
{Part of the cooperative-communication substrate. \textbf{Note:} venue was imprecisely listed as ``GEM Workshop''; corrected to ``GEM\textsuperscript{2} @ ACL 2025'' with page numbers during the audit.}

\entry{chen2025llmhanabi}
{Liang, F., Zheng, T., Chan, C., Yim, Y., Song, Y.\ (2025). LLM-Hanabi: Evaluating Multi-Agent Gameplays with Theory-of-Mind. arXiv:2510.04980.}
{https://arxiv.org/abs/2510.04980}
{Establishes that first-order theory-of-mind (interpreting another agent's intent) is the strongest predictor of cooperative success in LLM-Hanabi gameplay.}
{Critical for the ToM-erodes-with-language argument in Section 5: when first-order ToM degrades in lower-resource languages, both cooperation and deception suffer together.}

\entry{sun2025overcooked}
{Sun, H., Zhang, S., Niu, L., Ren, L., Xu, H., Fu, H., Zhao, F., Yuan, C., Wang, X.\ (2025). Collab-Overcooked: Benchmarking LLMs as Collaborative Agents. \emph{EMNLP 2025}.}
{https://arxiv.org/abs/2502.20073}
{Overcooked-based benchmark for evaluating LLM agents as collaborative real-time partners.}
{Part of the cooperative-communication substrate.}

\entry{chan2025xtom}
{Chan, C., Yim, Y., Zeng, H., Zou, Z., Cheng, X., Sun, Z., Deng, Z., Chung, K., Ao, Y., Fan, Y., Jiayang, C., Nie, E., Wong, G.Y., Schmid, H., Sch\"utze, H., See, S., Song, Y.\ (2025). XToM: Exploring the Multilingual Theory of Mind. arXiv:2506.02461.}
{https://arxiv.org/abs/2506.02461}
{Multilingual ToM benchmark documenting the ``pronounced dissonance'' between multilingual language understanding and ToM accuracy across languages on parallel items.}
{The empirical anchor for the multilingual-ToM-erosion finding in Section 5. \textbf{Note:} bib author list expanded from ``and others'' to full 17-author list during the audit.}

% =====================================================
\section{Open-world game environments}

\entry{wang2024voyager}
{Wang, G., Xie, Y., Jiang, Y., Mandlekar, A., Xiao, C., Zhu, Y., Fan, L., Anandkumar, A.\ (2024). Voyager: An Open-Ended Embodied Agent with LLMs. \emph{TMLR 2024}.}
{https://arxiv.org/abs/2305.16291}
{Open-ended Minecraft-playing LLM agent that incrementally acquires skills over long horizons.}
{Cited as the canonical open-world example in Section 2's substrate.}

\entry{paischer2024balrog}
{Paglieri, D., Cupia\l, B., Coward, S., et al.\ (2024). BALROG: Benchmarking Agentic LLM and VLM Reasoning On Games. \emph{ICLR 2025}.}
{https://arxiv.org/abs/2411.13543}
{Benchmark of agentic LLM and vision-language reasoning across challenging game environments.}
{Part of the open-world substrate.}

\entry{chen2025gvgai}
{Li, Y., Lin, C., Nasir, M.U., Bontrager, P., Liu, J., Togelius, J.\ (2025). GVGAI-LLM: Evaluating LLM Agents with Infinite Games. arXiv:2508.08501.}
{https://arxiv.org/abs/2508.08501}
{General Video Game AI extension for evaluating LLM agents on procedurally generated games.}
{Open-world substrate.}

\entry{wang2025tcgbench}
{AlRashed, S., Wang, J., Orabona, F.\ (2026). Cards Against Contamination: TCG-Bench for Difficulty-Scalable Multilingual LLM Reasoning. \emph{Findings of EACL 2026}.}
{https://aclanthology.org/2026.findings-eacl.353}
{Trading-card-game-style multilingual benchmark; documents a substantial English--Arabic performance gap at 32B parameters.}
{The empirical source for the ``scale alone does not correct the drift'' claim in Section 6.2 (cultural evolution).}

% =====================================================
\section{Dialogue games (clembench)}

\entry{chalamalasetti2023clembench}
{Chalamalasetti, K., G\"otze, J., Hakimov, S., Madureira, B., Sadler, P., Schlangen, D.\ (2023). clembench: Using Game Play to Evaluate Chat-Optimized LMs. \emph{EMNLP 2023}.}
{https://aclanthology.org/2023.emnlp-main.689/}
{Original clembench paper introducing dialogue-game-based evaluation of chat-tuned LLMs.}
{Founding paper of the clembench programme; cited as the dialogue-game family in Section 2.}

\entry{beyer2024clembench}
{Beyer, A., Chalamalasetti, K., Hakimov, S., Madureira, B., Sadler, P., Schlangen, D.\ (2024). clembench-2024: A Multilingual Benchmark for LLMs as Multi-Action Agents. arXiv:2405.20859.}
{https://arxiv.org/abs/2405.20859}
{2024 multilingual extension of clembench with format-adherence diagnostics across languages.}
{Source for the German/Italian late-layer realisation-failure observations cited in Section 4. \textbf{Note:} bib author list expanded from ``and others'' to full 6-author list during the audit; title hyphenation also corrected.}

\entry{schlangen2025clembench}
{Schlangen, D., Hakimov, S., Chalamalasetti, K., Jordan, J., Sadler, P.\ (2025). A Third Paradigm for LLM Evaluation. arXiv:2507.08491.}
{https://arxiv.org/abs/2507.08491}
{Position paper articulating dialogue-game-based evaluation as a third paradigm beyond supervised benchmarks and human ratings.}
{Cited as the methodological positioning of clembench. \textbf{Note:} bib author list expanded during the audit.}

\entry{hakimov2026pricethought}
{Hakimov, S., et al.\ (2026). The Price of Thought: A Multilingual Analysis of Reasoning, Performance, and Cost of Negotiation in LLMs.}
{https://arxiv.org/}
{University of Potsdam preprint analysing the cost-benefit trade-offs of multilingual reasoning in negotiation tasks.}
{Cited as part of the dialogue-games substrate in Section 2.}

% =====================================================
\section{Multi-agent benchmarks and failure analysis}

\entry{cemri2025mast}
{Cemri, M., Pan, M.Z., Yang, S., Agrawal, L.A., Chopra, B., Tiwari, R., Keutzer, K., Parameswaran, A., Klein, D., Ramchandran, K., Zaharia, M., Gonzalez, J.E., Stoica, I.\ (2025). Why Do Multi-Agent LLM Systems Fail? arXiv:2503.13657.}
{https://arxiv.org/abs/2503.13657}
{MAST taxonomy: 14 failure modes derived from 1{,}600+ execution traces across 7 multi-agent frameworks; categorises failures into Specification (41.77\%), Inter-Agent Misalignment (36.94\%), Task Verification (21.30\%).}
{The empirical anchor for the ``intent-action gap'' discussion in Section 5; the categorisation supports the cognitive framing of multi-agent failures.}

\entry{nie2025multiagentbench}
{Zhu, K., Du, H., Hong, Z., Yang, X., Guo, S., Wang, Z., Wang, Z., Qian, C., Tang, X., Ji, H., You, J.\ (2025). MultiAgentBench: Evaluating the Collaboration and Competition of LLM Agents. \emph{ACL 2025}.}
{https://arxiv.org/abs/2503.01935}
{Comprehensive benchmark for evaluating LLM-agent collaboration and competition across diverse tasks.}
{Cited in Section 2's substrate and as one of the interactive multi-agent benchmarks to be paired with safety benchmarks in our governance recommendations.}

\entry{xie2026m3bench}
{Xie, S., Shi, Z., Shen, H., Huang, G., Ma, Y., Jing, X.\ (2026). M3-BENCH: Process-Aware Evaluation of LLM Agents' Social Behaviors in Mixed-Motive Games. arXiv:2601.08462.}
{https://arxiv.org/abs/2601.08462}
{Process-aware multi-agent benchmark with a commitment-reliability metric showing promise-keeping drops in mixed-motive multilingual games.}
{Source for the commitment-reliability claim in Section 5 and for the broader cognitive-failures framing.}

\entry{sun2025gametheoryllm}
{Sun, H., Wu, Y., Wang, P., Chen, W., Cheng, Y., Deng, X., Chu, X.\ (2025). Game Theory Meets LLMs: A Systematic Survey. \emph{IJCAI 2025}.}
{https://www.ijcai.org/}
{IJCAI survey-track paper systematising the intersection of game theory and LLM research.}
{Cited as the standing survey of the field in Section 1. \textbf{Note:} bib entry was missing 3 authors (Wang, Chen, Deng); corrected during the audit.}

% =====================================================
\section{Pretraining-corpus audits}

\entry{gao2020pile}
{Gao, L., Biderman, S., Black, S., Golding, L., Hoppe, T., Foster, C., Phang, J., He, H., Thite, A., Nabeshima, N., Presser, S., Leahy, C.\ (2020). The Pile: An 800GB Dataset of Diverse Text. arXiv:2101.00027.}
{https://arxiv.org/abs/2101.00027}
{Documents the construction and contents of the Pile, an 800GB diverse-text pretraining corpus.}
{Cited in P1 as the open-data corpus that supplies ground truth for measuring cooperative-discourse density.}

\entry{luccioni2021what}
{Luccioni, A.S., Viviano, J.D.\ (2021). What's in the Box? An Analysis of Undesirable Content in the Common Crawl Corpus. \emph{ACL 2021}.}
{https://aclanthology.org/2021.acl-short.24/}
{Audits Common Crawl for undesirable content (toxicity, bias) and provides methodology for corpus inspection.}
{Cited together with the Pile paper as the methodology for P1's corpus-audit operationalisation.}

% =====================================================
\section{Other supporting citations}

\entry{whorf1956language}
{Whorf, B.L.\ (1956). \emph{Language, Thought, and Reality}. MIT Press.}
{https://mitpress.mit.edu/9780262730068/language-thought-and-reality/}
{The collected writings that established the strong linguistic-relativity hypothesis Whorf is named after.}
{Cited only in passing as the historical anchor of the term ``Sapir-Whorf'' which we explicitly weaken.}

\entry{ponti2020xcopa}
{Ponti, E.M., Glava\v{s}, G., Majewska, O., Liu, Q., Vuli\'c, I., Korhonen, A.\ (2020). XCOPA: A Multilingual Dataset for Causal Commonsense Reasoning. \emph{EMNLP 2020}.}
{https://aclanthology.org/2020.emnlp-main.185/}
{Multilingual extension of COPA for causal commonsense reasoning across 11 typologically diverse languages.}
{Background reference for multilingual evaluation; cited in early drafts.}

% =====================================================
\section*{End notes}

\textbf{Coverage.} This document covers all entries cited at least once in the published version of \emph{Babel's Machines} as of the Round 2 revision (2026-04-27). A handful of additional entries appear in \texttt{refs.bib} but are not currently cited inline; those are not described here.

\textbf{If the paper is revised further.} New citations should be web-verified at insertion time (per the user's saved feedback memory: ``verify every bib entry against arXiv/ACL Anthology before citing''), and this companion document updated.

\textbf{Audit history.}
\begin{itemize}
  \item Round 1 audit (2026-04-27): 73 cited keys checked. 12 entries needed substantive fixes (wrong title, wrong venue, missing authors, wrong year). 60 verified clean. 1 cannot fully verify (\texttt{repello2024multilingualguardrails}, blog post). All fixes applied to \texttt{refs.bib} on the same day.
  \item Round 2 additions (2026-04-27): 7 new bib entries (5 interpretability precedents from reviewer + 2 ToM critique papers) all web-verified before insertion.
  \item Round 1 reviewer additions: 9 new entries (3 Round 1 reviewer-cited + 3 multilingual rep classics + 3 language-prioritised models) all web-verified before insertion.
\end{itemize}

\end{document}
